\section{Dữ liệu tensor stream và bài toán dự báo}
\label{sec:problem}

\subsection{Từ bảng đến tensor}

Một chuỗi thời gian thông thường có dạng $x_1,x_2,\ldots,x_T$: một giá trị tại mỗi thời điểm.
Dữ liệu của D-Tracker có thêm hai chiều. Với $k$ từ khóa, $l$ địa điểm và $t$ thời điểm, mỗi
lát cắt thời gian là một \emph{ma trận} $X_t\in\mathbb{R}^{k\times l}$, trong đó
$x_{ij}(t)$ là mức quan tâm của từ khóa $i$ tại địa điểm $j$ ở thời điểm $t$. Ghép các ma trận
lại, ta có tensor bậc ba
\begin{equation}
  \mathcal{X}=\{X_1,\ldots,X_{t_c}\}\in\mathbb{N}^{k\times l\times t_c},
  \label{eq:stream}
\end{equation}
với $t_c$ là thời điểm hiện tại. Paper hiển thị tensor thực nghiệm theo thứ tự thời gian, từ khóa,
địa điểm; khi viết công thức, tác giả dùng các chiều $k,l,L_c$. Hai cách viết mô tả cùng một đối
tượng, chỉ khác thứ tự trục.

\begin{center}
\small
\begin{tabular}{@{}lll@{}}
  \toprule
  Ký hiệu & Ý nghĩa & Ví dụ \\
  \midrule
  $k$ & số từ khóa & iPhone, Galaxy, iPad \\
  $l$ & số địa điểm & Hoa Kỳ, Việt Nam, Nhật Bản \\
  $t_c$ & thời điểm hiện tại & tuần thứ 676 \\
  $L_c$ & độ dài cửa sổ hiện tại & 104 tuần (hai năm) \\
  $L_f$ & tầm dự báo & 13, 26 hoặc 39 tuần \\
  \bottomrule
\end{tabular}
\end{center}

\subsection{Stream: dữ liệu không có điểm kết thúc}

Trong ứng dụng thật, mỗi tuần ta nhận thêm một ma trận $X_{t_c+1}$. Không thể giữ toàn bộ lịch
sử mãi mãi, và mô hình cũ có thể không còn đúng. Paper định nghĩa \emph{cửa sổ hiện tại} bằng
$L_c$ lát cắt gần nhất. Bài in công thức với hai đầu mút
$\{X_t\}_{t=t_c-L_c}^{t_c}$; nếu hiểu cả hai đầu mút theo nghĩa bao gồm thì đó là $L_c+1$ lát cắt,
nên tài liệu dùng dạng nhất quán với độ dài đã nêu:
\begin{equation}
  \mathcal{X}^{c}=\{X_{t_c-L_c+1},\ldots,X_{t_c}\},
  \label{eq:window}
\end{equation}
gồm đúng $L_c$ lát cắt. Đây là điểm rất quan trọng: D-Tracker không học lại từ toàn bộ
stream ở mỗi bước. Nó ước lượng lại trên cửa sổ hiện tại, rồi dùng cơ chế cập nhật để quyết định
có thêm một mô hình mới vào bộ mô hình hay không.

\begin{center}
\begin{tikzpicture}[
  >=Latex,
  box/.style={draw=DTNavy, fill=DTMist, rounded corners=1.5pt,
              minimum height=8mm, align=center, font=\sffamily\small},
  obs/.style={draw=DTTeal, fill=white, minimum width=9mm, minimum height=8mm,
              align=center, font=\small},
  future/.style={draw=DTClaret, dashed, fill=DTWarm, minimum width=12mm,
                 minimum height=8mm, align=center, font=\small}
]
  \node[font=\sffamily\bfseries\color{DTNavy}] at (0,1.15) {Quá khứ};
  \node[font=\sffamily\bfseries\color{DTTeal}] at (5.1,1.15) {Cửa sổ hiện tại $L_c$};
  \node[font=\sffamily\bfseries\color{DTClaret}] at (10.3,1.15) {Tương lai $L_f$};
  \draw[->, thick, DTNavy] (-0.5,0) -- (12,0);
  \foreach \x in {0,...,4} {
    \node[obs] (h\x) at (\x*1.05,0) {};
  }
  \foreach \x in {0,...,2} {
    \node[obs, draw=DTTeal, fill=DTTeal!12] (c\x) at (5.1+\x*1.05,0) {};
  }
  \node[future] (f1) at (10.05,0) {$\cdots$};
  \node[future] (f2) at (11.35,0) {$X_{t_c+1}$};
  \draw[decorate,decoration={brace,amplitude=4pt},DTTeal]
    (4.6,0.55) -- (7.65,0.55) node[midway,above=3pt,font=\small] {$L_c$};
  \draw[decorate,decoration={brace,amplitude=4pt},DTClaret]
    (9.45,0.55) -- (11.85,0.55) node[midway,above=3pt,font=\small] {$L_f$};
  \node[box, below=8mm of c1, xshift=6mm] {dùng để ước lượng mô hình hiện tại};
  \node[box, below=8mm of f2] {cần dự báo};
\end{tikzpicture}
\end{center}

\subsection{Bài toán chính thức}

Paper phát biểu bài toán \emph{stream forecasting} như sau: cho cửa sổ $\mathcal{X}^{c}$ dài
$L_c$, hãy tìm các động học thay đổi theo thời gian gồm trend, mùa vụ và lan truyền, rồi dự báo
giá trị tại tầm $L_f$. Đầu ra là ma trận tại thời điểm đích
$\widehat{X}_{t_c+L_f}\in\mathbb{R}_{+}^{k\times l}$; khi đánh giá nhiều horizon, paper chạy từng
tầm riêng và không nói rằng mô hình phải sinh mọi ma trận trung gian.

\begin{keypoint}
\textbf{Đầu vào:} $L_c$ ma trận quan sát gần nhất, mỗi ma trận kích thước $k\times l$.\\
\textbf{Đầu ra:} dự báo tại thời điểm $t_c+L_f$, cùng một mô hình có các thành phần trend,
seasonal, outlier và lan truyền có thể đọc được.\\
\textbf{Ràng buộc:} không giữ toàn bộ stream; mỗi bước chỉ được xử lý trong thời gian không
phụ thuộc vào tổng chiều dài đã quan sát.
\end{keypoint}

Trong ví dụ ở Chương~\ref{sec:orientation}, $k=2$ (iPhone, Galaxy), $l=2$ (Hoa Kỳ, Việt Nam),
$L_c=3$ (ba tuần đã có), và $L_f=1$ nếu ta muốn dự báo tuần kế tiếp. Một mô hình tốt không chỉ
khớp ba tuần đó; nó phải tách được xu hướng của từng cặp và biểu diễn được việc mức quan tâm
có thể lan giữa hai địa điểm.

\subsection{Vì sao không dự báo từng chuỗi riêng lẻ?}

Ta có thể chạy một mô hình chuỗi thời gian cho từng cặp $(i,j)$. Cách đó bỏ qua hai điều:
\begin{itemize}
  \item các địa điểm và từ khóa chia sẻ những \emph{nhóm động học} tiềm ẩn; và
  \item dữ liệu có cấu trúc đa tuyến tính, nên một hệ số nhỏ có thể mô tả đồng thời nhiều cặp.
\end{itemize}
D-Tracker tìm một lõi nhỏ $\mathcal{W}^{(\mathrm{core})}$ và các hệ số địa điểm/từ khóa để tái tạo
toàn bộ tensor. Nhờ đó mô hình vừa giảm số tham số vừa tạo ra các nhóm có thể diễn giải. Cái giá
là phải giải một bài toán tối ưu khó hơn và phải kiểm tra giả định rằng cấu trúc nhóm đó thật sự
hữu ích.

\begin{plainwords}
Đừng nghĩ ``tensor'' là một khối dữ liệu bí ẩn. Hãy nghĩ đến một chồng bảng: mỗi tuần là một
bảng từ khóa $\times$ địa điểm. D-Tracker tìm một vài kiểu bảng chung, rồi trộn chúng với trọng
số khác nhau cho từng từ khóa và địa điểm.
\end{plainwords}

\begin{commonmistake}
Có hai nhầm lẫn thường gặp. Thứ nhất, tensor bậc ba không phải ``ba chiều không gian''; ba trục
ở đây là các thuộc tính dữ liệu. Thứ hai, mô hình không dự báo toàn bộ stream vô hạn cùng lúc:
nó chỉ dự báo $L_f$ bước từ cửa sổ hiện tại, rồi lặp lại khi có dữ liệu mới.
\end{commonmistake}

\subsection{Ba câu hỏi mà mô hình phải trả lời}

Paper đánh giá mô hình theo ba câu hỏi, và ta sẽ dùng chúng như khung đọc:
\begin{enumerate}
  \item \textbf{Accuracy:} dự báo tương lai chính xác đến đâu? (MAE, RMSE)
  \item \textbf{Scalability:} thời gian chạy thay đổi thế nào khi stream dài hơn?
  \item \textbf{Effectiveness:} mô hình có tách được trend, mùa vụ và lan truyền theo cách
        diễn giải được không?
\end{enumerate}
Ba câu hỏi này độc lập: một mô hình có thể khớp tốt nhưng không giải thích được, hoặc giải thích
đẹp nhưng dự báo kém. Toàn bộ thiết kế của D-Tracker là một nỗ lực cân bằng cả ba.

\begin{selfcheck}
\begin{enumerate}
  \item Với tensor $k=8,l=50,L_c=104$, mỗi lát cắt thời gian có bao nhiêu phần tử?
  \item Phân biệt $t_c$, $L_c$ và $L_f$ trong một câu.
  \item Tại sao chỉ dùng cửa sổ gần nhất lại giúp mô hình thích nghi với pattern thay đổi?
\end{enumerate}
\end{selfcheck}
